% ===================================================================
% Appendix to Section 11 (Contexts, idealization, and checking physics)
% Source: research/theory/T3-contexts-idealization-export.md, as repaired after
%         verification (see its "Verification log": A1-A18, B1-B13, C1-C11);
%         research/theory/T2-coherence-as-negative-data.md (Prop 7.1, proved in the
%         main text).
% ===================================================================
\providecommand{\fM}{\mathfrak{M}}
\providecommand{\fF}{\mathcal{F}}
\providecommand{\fG}{\mathcal{G}}
\providecommand{\fK}{\mathcal{K}}
\providecommand{\tset}[1]{\lVert #1\rVert}
\providecommand{\cder}{\Vdash}
\providecommand{\dom}{\operatorname{dom}}
\providecommand{\osc}{\operatorname{osc}}
\providecommand{\Anc}{\mathrm{Anc}}
\providecommand{\Asm}{\mathrm{Asm}}
\providecommand{\Ax}{\mathrm{Ax}}

\section{Proofs for Section~\ref{sec:physics}}\label{app:physics}

The following results are proved completely in the main text:
\cref{prop:physics:los,prop:physics:sup,thm:physics:eternalism,prop:physics:misdesignation,thm:physics:soundness,prop:physics:continuity,lem:physics:invisible,prop:physics:radius,thm:physics:coherencetol,thm:physics:judgment,prop:physics:gricean}.
The minimal-repair drift stated after \cref{thm:physics:coherencetol} (T3 Prop~4.2) is
immediate: enlarging a tolerance enlarges an interval, so coherence is preserved upward.
\Cref{prop:physics:laymon} and the numerical
claims marked \status{computed} are computations, reproduced by the scripts in
\texttt{research/theory/T3-checks/}; they are not proved here. Everything else is
proved below. Throughout, $w$ is a world, $\dom=\dom\fM_w$, and
$\delta(F\times G):=\{\lambda[D{:=}\mu]:\lambda\in F,\ \mu\in G\}$.

% -------------------------------------------------------------------
\subsection{Preliminaries on the filters of Definition~\ref{def:physics:filter}}

\begin{lemma}[Filter bases]\label{lem:physics:base}
\textup{(a)} For $c=\DEF(D,\lambda^\circ,\fG)$ with parent $p$, the family
$\mathcal B_c=\{\delta(F\times G)\cap\dom:F\in\fF_p(w),\ G\in\fG\}$ is a filter base:
$X\in\fF_c(w)$ iff $X\supseteq B$ for some $B\in\mathcal B_c$, and $c$ is proper at $w$
iff every member of $\mathcal B_c$ is nonempty.
\textup{(b)} For $c=\SUP(A)$ with parent $p$, $X\in\fF_c(w)$ iff $X\supseteq F\cap\tset{A}_w$
for some $F\in\fF_p(w)$.
\textup{(c)} If the parent filter is principal at a point $\lambda$, then in (a) one may
take $F=\{\lambda\}$: $X\in\fF_c(w)$ iff $X\supseteq\{\lambda[D{:=}\mu]:\mu\in G\}\cap\dom$
for some $G\in\fG$.
\end{lemma}
\begin{proof}
(a) $\delta(F_1\times G_1)\cap\delta(F_2\times G_2)\supseteq\delta((F_1\cap F_2)\times(G_1\cap G_2))$,
and $F_1\cap F_2\in\fF_p(w)$, $G_1\cap G_2\in\fG$. So $\mathcal B_c$ is directed downward,
and the filter it generates consists of the supersets (in $\dom$) of its members. It
contains $\emptyset$ iff some member is empty. (b) Likewise $\{F\cap\tset{A}:F\in\fF_p(w)\}$
is closed under finite intersections. (c) If $\fF_p(w)$ is principal at $\lambda$, every
$F\in\fF_p(w)$ contains $\lambda$, so $\delta(\{\lambda\}\times G)\subseteq\delta(F\times G)$,
and $\{\lambda\}\in\fF_p(w)$.
\end{proof}

% -------------------------------------------------------------------
\subsection{Proofs for \S\ref{sec:physics:semantics}}

\begin{proof}[Proof of \cref{prop:physics:idealized}]
(a) The root filter is principal at $\lambda^*$, where $p_{\rm atm}=1.0\times10^5$. The
frame's laws do not mention phase change, so the model exists at
$\lambda^*[p_{\rm atm}{:=}0]$, and by \cref{lem:physics:base}(c) with the principal path
the child filter is principal there; hence $c\models p_{\rm atm}=0$. Both filters are
principal, hence proper, and by \cref{prop:physics:los}(b) each context's theory is
consistent.

(b) Under limit semantics the only generator is $\{\lambda^*[m_r{:=}0]\}\cap\dom=\emptyset$,
so the limit child is improper. Under germ semantics the generators are
$\{\lambda^*[m_r{:=}\mu]:0<|\mu|<\eta\}\cap\dom=\{\lambda^*[m_r{:=}\mu]:0<\mu<\eta\}$, all
nonempty, so the germ child is proper (\cref{lem:physics:base}(a)). In the model at
$m_r=\mu>0$, $a=\Phi/\mu$, so $|a|>K$ whenever $\mu<|\Phi|/\max(K,1)$; the generator with
$\eta=|\Phi|/\max(K,1)$ lies inside $\tset{|a|>K}$.

(c) Let $E=\{\mu:\lambda^*[D{:=}\mu]\in\dom\}$. By \cref{lem:physics:base}(c),
$c_{\rm germ}\models Q\in I$ iff there is $G\in\fG$ with $q(\mu)\in I$ for all $\mu\in G\cap E$.
``For every open $I\ni q(\lambda^\circ)$ there is $G\in\fG$ with $q(G\cap E)\subseteq I$'' is
the definition of $q(\mu)\to q(\lambda^\circ)$ along the trace of $\fG$ on $E$. (Intervals
with endpoints named in $L$, e.g.\ rational ones, suffice for convergence.) The limit
child is principal at $\lambda^*[D{:=}\lambda^\circ]$, so its value of $Q$ is
$q(\lambda^\circ)$.
\end{proof}

\begin{proof}[Proof of \cref{lem:physics:import}]
(a) Let $F=\tset{\varphi}_w\in\fF_p(w)$ and $G\in\fG$. Every element of
$\delta(F\times G)\cap\dom$ has the form $\lambda[D{:=}\mu]\in\dom$ with $\lambda\in F$; it
differs from $\lambda$ only in coordinates in $D$, so by $D$-stability it satisfies
$\varphi$. Thus this generator lies inside $\tset{\varphi}_w$, and
$\tset{\varphi}_w\in\fF_c(w)$. Laws are $D$-stable because they hold on all of $\dom$.
An arithmetic formula in the constants $p_i$, $i\notin D$, is evaluated in the standard
real sort with $p_i$ interpreted as $\lambda_i$, so its truth value depends only on
$(\lambda_i)_{i\notin D}$, which points differing only on $D$ share.

(b) Let $\lambda,\lambda'\in\dom$ differ only on $D$, with $\fM_w(\lambda)\models\varphi$.
Take the actual point $\lambda^*=\lambda$ and the ideal value $\lambda^\circ=\lambda'_D$.
Then $@\models\varphi$, and the limit child's only generator is
$\{\lambda[D{:=}\lambda'_D]\}\cap\dom=\{\lambda'\}$, so its filter is principal at $\lambda'$
and $c\models\varphi$ says $\fM_w(\lambda')\models\varphi$. Exchanging the roles of
$\lambda$ and $\lambda'$ gives the converse.

(c) Every generator $\{\lambda[k{:=}0]:\lambda\in F\}\cap\dom$ of the child lies inside
$\tset{k=0}$, so $c\models k=0$ and, by \cref{prop:physics:los}(a), $c\models k\le1$, for
every parent filter and every world (if $c$ is improper it proves everything). So an
import of $\varphi$ into $c$ never leads from a truth to a falsehood. But if $\dom$ contains
$\lambda$ and $\lambda'$ differing only in $k$, with $k=0.5$ and $k=2$ respectively, then
$\fM_w(\lambda)\models\varphi$ and $\fM_w(\lambda')\not\models\varphi$.
\end{proof}

% -------------------------------------------------------------------
\subsection{Proofs for \S\ref{sec:physics:coherence}: realizability}

$\AUX$ contexts are identified with their parents throughout. For a set $M\subseteq\fK$ and a
sentence $A$, $\tset{A}$ in (R1) denotes $\{N\in\fK:N\models A\}$. Since $\SUP$ contexts
have only $\SUP$ descendants, the parent of a $\DEF$ context is never a $\SUP$ context.

\begin{proof}[Proof of \cref{thm:physics:realizability}(a)]
\emph{Observation 1 (collapse).} By (R1) and induction along $\SUP$ chains,
$M_c=M_b\cap\bigcap_{A\in\Asm(c)}\tset{A}$ for every $\SUP$ context $c$ with base $b$. Hence
$(c,\Gamma,\varphi)$ satisfies (R3) at $M_c$ iff every $N\in M_b$ with
$N\models\bigwedge\Asm(c)\wedge\bigwedge\Gamma$ satisfies $\varphi$, i.e.\ iff
$(b,\Gamma\cup\Asm(c),\varphi)$ satisfies (R3) at $M_b$. So, given (R1), the judgments of $J$
hold iff $M_b\subseteq\{N:N\models\Th^*(b)\}$ for every non-$\SUP$ context $b$.

\emph{Observation 2 (anchoring lemma).} In any realization, $M_c\neq\emptyset$ for every
$c\in\Anc$. We induct along the inductive definition of $\Anc$. The root is nonempty by
(R0) and certified contexts by (R5). Suppose $c$ enters through an informative use
$(c,\beta,t)$ with $M_{\pi c}\neq\emptyset$, and suppose $M_c=\emptyset$. Then
$M_c\models\phi_\beta(v)$ for every $v\in V$ (vacuously), so (R4) gives
$M_{\pi c}\models\sigma_\beta\to\varepsilon_\beta(v)$ for all $v$. Since
$(\pi c,\emptyset,\sigma_\beta)\in J$, (R3) gives $M_{\pi c}\models\sigma_\beta$. Any
$N\in M_{\pi c}$ then satisfies $\{\sigma_\beta\}\cup\{\varepsilon_\beta(v):v\in V_0\}$,
contradicting informativeness, since $N\in\fK$.

($\Rightarrow$) For $c\in\Anc$, $M_c\neq\emptyset$ by Observation~2, and any $N\in M_c$
satisfies $\Th^*(c)$ by Observation~1.

($\Leftarrow$) For each $c\in\Anc$ choose $N_c\in\fK$ with $N_c\models\Th^*(c)$, and put
$M_c=\{N_c\}$. Put $M_c=\emptyset$ for every non-$\SUP$ context outside $\Anc$ (these are
$\DEF$ contexts, since $@\in\Anc$), and define $\SUP$ contexts by (R1). Then (R0), (R1),
(R2) and (R5) hold by construction, since $@\in\Anc$ and $C\subseteq\Anc$. (R3) holds by
Observation~1: at anchored $b$ because $N_b\models\Th^*(b)$, and at non-anchored $b$
vacuously. For (R4), take a use $(c,\beta,t)$. If $\pi c\notin\Anc$, then
$M_{\pi c}=\emptyset$ and (R4) is vacuous. If $\pi c\in\Anc$, then $c\in\Anc$ because
$\beta$ is informative, so $M_c=\{N_c\}$. The judgment $(c,\emptyset,Q=\bar t)$ lies in
$J^*$ unchanged ($c$ is not a $\SUP$ context), so $Q^{N_c}=t$; by standard names,
$M_c\models Q=\bar v$ only for $v=t$. Finally $(\pi c,\emptyset,\varepsilon_\beta(t))\in J$
gives $N_{\pi c}\models\varepsilon_\beta(t)$, so $M_{\pi c}\models\sigma_\beta\to\varepsilon_\beta(t)$.
\end{proof}

\begin{proof}[Proof of \cref{thm:physics:realizability}(b)]
Fix a frame realization at $w$.

\emph{Filter anchoring lemma: every $c\in\Anc$ is proper at $w$.} The root is proper, and
certified contexts are proper by assumption. Suppose $c$ enters through an informative use
$(c,\beta,t)$ with $\pi c$ proper, and suppose $c$ is improper. Then $c\models_w\phi_\beta(v)$
for every $v$, so bridge soundness gives $\pi c\models_w\sigma_\beta\to\varepsilon_\beta(v)$
for all $v$. The judgment $(\pi c,\emptyset,\sigma_\beta)$ holds, so
$\pi c\models_w\varepsilon_\beta(v)$ for all $v$ (\cref{prop:physics:los}(a)). By
informativeness, $\{\sigma_\beta\}\cup\{\varepsilon_\beta(v):v\in V_0\}$ has no model in $\fK$
for some finite $V_0$; every $\fM_w(\lambda)$ is in $\fK$, so the truth set of the finite
conjunction is $\emptyset$. That set lies in $\fF_{\pi c}(w)$, contradicting properness.

\emph{Conclusion.} Let $c\in\Anc$. For a $\SUP$ descendant
$c'=\SUP(A_k)$ of $\dots$ of $\SUP(A_1)$ of $c$, iterating \cref{prop:physics:sup} and
using \cref{prop:physics:los}(a) shows that $c'\models_w\chi$ iff
$c\models_w\bigwedge_iA_i\to\chi$. Hence each judgment $(c',\Gamma,\varphi)$ holds iff
$c\models_w\bigwedge\Asm(c')\wedge\bigwedge\Gamma\to\varphi$, and so $\Th^*(c)\subseteq\Th_w(c)$.
Since $\Th^*(c)$ is finite, the intersection of the truth sets of its members lies in
$\fF_c(w)$, and it is nonempty because $c$ is proper. Any $\lambda$ in it gives a model
$\fM_w(\lambda)\in\fK$ of $\Th^*(c)$. No compactness is used. The argument does not
depend on whether paths are principal or germs.
\end{proof}

\begin{proof}[Proof of \cref{thm:physics:realizability}(c)]
\emph{CE1.} $C=\{c\}$, so $\Anc=\{@,c\}$; $\Th^*(@)=\{p_1=1\}$ and $\Th^*(c)=\{p_1=0\}$ are
each satisfiable in $\fK$ (the constant $p_1$ may denote any real), so the criterion holds.
In a frame realization, the root judgment forces $\lambda^*_1=1$. The certified $c$ is
proper, so $\lambda^*[2{:=}0]\in\dom$ and $\fF_c$ is principal at that point, whose first
coordinate is $1$; hence $c\models p_1=1$, and $c\not\models p_1=0$ since $\Th_w(c)$ is
consistent. \emph{CE2.} Two siblings with the same declaration have the same generators,
hence the same filter. Both are certified, hence proper; the judgments $Q=1$ and $Q=2$
would put $\tset{Q=1}\cap\tset{Q=2}=\emptyset$ in this proper filter. The criterion holds
because each of $\{Q=1\}$, $\{Q=2\}$ is satisfiable. Both counterexamples are also
brute-forced over all frames on $\Lambda=\{0,1\}^2$ (\texttt{realizability\_frame.py}).
\end{proof}

The exact frame criterion under limit semantics, summarized in the main text after
\cref{thm:physics:realizability}, is the following.

\begin{proposition}[Frame realizability under limit semantics]\src{T3 Prop 2.7, added after verification}\ \status{brute-force checked}\label{prop:physics:framerealizability}
Assume every $\DEF$ context in $J$ uses limit semantics, every used bridge is an informative
value-export bridge, and $\SUP$ contexts have only $\SUP$ descendants. For a candidate actual
point $\lambda$, define the points of the non-$\SUP$ contexts by $\lambda_@=\lambda$ and
$\lambda_c=\lambda_{\pi c}[D_c{:=}\lambda^\circ_c]$. Let $S(\lambda)$ be the least set of
non-$\SUP$ contexts that contains $@$ and $C$ and is closed under (i) \emph{parents}:
$c\in S\Rightarrow\pi c\in S$; (ii) \emph{informative bridges}: $(c,\beta,t)\in U$ and
$\pi c\in S\Rightarrow c\in S$; (iii) \emph{coincidence}: $\pi c\in S$ and
$\lambda_c=\lambda_{c'}$ for some $c'\in S\Rightarrow c\in S$. Then $J$ is frame-realizable iff
there is $\lambda\in\Lambda$ such that every point $x\in\{\lambda_c:c\in S(\lambda)\}$ lies in
$\Lambda$ and $\bigcup\{\Th^*(c):c\in S(\lambda),\ \lambda_c=x\}$ has a model $N\in\fK$ with
$p_i^N=x_i$ for all $i$.
\end{proposition}

The criterion is non-local in two ways: it quantifies over the actual parameter values, and
it conjoins the theories of all forced-proper contexts that sit at one parameter point. CE1
fails it because $c$'s point shares $p_1$ with the root, and CE2 because the two siblings
share a point. \texttt{realizability\_frame.py} compares it with brute-force frame
realizability over all frames on $\Lambda=\{0,1\}^2$ (up to three $\DEF$ and two $\SUP$
contexts, $0$--$2$ bridge uses per context): $8{,}000$ instances, $0$ mismatches.

\begin{proof}[Proof of \cref{prop:physics:framerealizability}]
\emph{Shape of the filters.} Under limit semantics every non-$\SUP$ context's filter is
improper or principal at its point. Indeed $\fF_@$ is principal at $\lambda^*=\lambda_@$. If
$\fF_{\pi c}$ is principal at $\lambda_{\pi c}$, \cref{lem:physics:base}(c) with $G=\{\lambda^\circ_c\}$
shows that $\fF_c$ is generated by $\{\lambda_c\}\cap\dom$: principal at $\lambda_c$ if
$\lambda_c\in\dom$, improper otherwise. If $\fF_{\pi c}$ is improper, so is $\fF_c$, since
$\delta(\emptyset\times G)=\emptyset$. A $\SUP$ context below a base principal at $x$ is
principal at $x$ or improper (\cref{lem:physics:base}(b)), and its judgments collapse into
the base by \cref{prop:physics:sup}.

($\Rightarrow$) Let $(\fM,\lambda^*)$ be a frame realization and let $P$ be the set of
proper non-$\SUP$ contexts. $P$ contains $@$ and $C$. It is closed under (i), since an
improper parent has an improper child; under (ii), by the filter anchoring lemma in the
proof of \cref{thm:physics:realizability}(b); and under (iii): if $\pi c\in P$ and
$\lambda_c=\lambda_{c'}$ with $c'\in P$, then $\lambda_{c'}\in\dom$ (a proper context is
principal at its point), so $\fF_c$ is principal at $\lambda_c\in\dom$ and $c$ is proper.
Hence $S(\lambda^*)\subseteq P$. For each point $x=\lambda_c$ with $c\in S(\lambda^*)$ we have
$x\in\dom\subseteq\Lambda$; every $c'\in S(\lambda^*)$ at $x$ is principal at $x$, so
$N=\fM(x)$ satisfies all their collapsed theories, and $p_i^N=x_i$.

($\Leftarrow$) Fix such a $\lambda$ and put $S=S(\lambda)$, $\lambda^*=\lambda$,
$\dom=\{\lambda_c:c\in S\}\subseteq\Lambda$, and let $\fM(x)$ be the given model at $x$; this
is a frame, since $p_i^{\fM(x)}=x_i$, and $\lambda^*=\lambda_@\in\dom$.
\emph{Contexts in $S$ are proper:} all their ancestors are in $S$ by (i), so all the
ancestors' points are in $\dom$, and by induction each is principal at its point.
\emph{Non-$\SUP$ contexts outside $S$ are improper:} by induction on depth; if
$\pi c\notin S$, then $\pi c$ is improper and so is $c$; if $\pi c\in S$, then (iii) gives
$\lambda_c\neq\lambda_{c'}$ for all $c'\in S$, i.e.\ $\lambda_c\notin\dom$.
\emph{Judgments} hold at a proper $c\in S$, because $\fM(\lambda_c)$ satisfies
$\Th^*(c)$, which by the collapse covers the judgments of $c$ and of its $\SUP$
descendants; at improper contexts they hold vacuously. \emph{Certified contexts} lie in
$S$, hence are proper. \emph{Bridge uses} $(c,\beta,t)$: if $\pi c$ is improper, soundness
is vacuous. Otherwise $\pi c\in S$, so $c\in S$ by (ii), and $c$ is principal at
$\lambda_c$. Then $c\models Q=\bar v$ iff $Q^{\fM(\lambda_c)}=v$, which by the judgment
$(c,\emptyset,Q=\bar t)$ and standard names happens only for $v=t$; and
$\pi c\models\varepsilon_\beta(t)$ by the judgment $(\pi c,\emptyset,\varepsilon_\beta(t))$.
\end{proof}

\begin{proof}[Proof of \cref{cor:physics:immunity}]
(i) In the frame-free sense, by the collapse a judgment $(c,\Gamma,\bot)$ at a $\SUP$
context $c$ with base $b$ becomes $\neg(\bigwedge\Asm(c)\wedge\bigwedge\Gamma)\in\Th^*(b)$.
In the frame semantics the same follows from \cref{prop:physics:sup}, iterated as in the
proof of \cref{thm:physics:realizability}(b). Adding this sentence to $\Th^*(b)$ makes a
satisfiable theory unsatisfiable only if the theory already entails
$\bigwedge\Asm(c)\wedge\bigwedge\Gamma$.

(ii) Let $c$ be a $\DEF$ context with $(c,\emptyset,\bot)\in J$; then $\bot\in\Th^*(c)$. If
$c\in\Anc$, $\Th^*(c)$ is unsatisfiable, so $J$ is not frame-free realizable by
\cref{thm:physics:realizability}(a), and not frame-realizable by (b). If $c\notin\Anc$:
$\Anc$ depends only on $C$ and $U$, and removing $(c,\emptyset,\bot)$ does not affect $U$
(bridge uses involve only judgments of the forms $Q=\bar t$, $\sigma_\beta$,
$\varepsilon_\beta(t)$; we assume $\bot$ is none of these). So by (a), $J$ is frame-free
realizable iff $J\setminus\{(c,\emptyset,\bot)\}$ is: the judgment imposes no constraint. In
the frame semantics this ``only if'' direction can fail. Example: $J=\{(@,\emptyset,p_1=0),
(c,\emptyset,\bot)\}$ with $c$ an uncertified limit child deforming $D=\{1\}$ toward $0$ and
no bridges. Then $c\notin\Anc$ and $J$ is frame-free realizable, but the root forces
$\lambda^*_1=0$, so $c$ sits at $\lambda^*\in\dom$, is proper, and cannot satisfy $\bot$.

(iii) Let $N\in\fK$ satisfy the stripped set. Put $M_c=\{N\}$ at every non-$\SUP$ context
and define $\SUP$ contexts by (R1). (R0), (R2), (R5) hold. (R3): each judgment
$(c,\Gamma,\varphi)$ has $N\models\bigwedge\Gamma\to\varphi$, and $M_c\subseteq\{N\}$. (R4): for
a use $(c,\beta,t)$, $N\models Q=\bar t$, so $M_c\models Q=\bar v$ only for $v=t$, and
$N\models\varepsilon_\beta(t)$; here $\pi c$ is not a $\SUP$ context, so $M_{\pi c}=\{N\}$. The
same construction works if $N$ satisfies $\bigcup_c\Th^*(c)$ (the collapsed variant), using
Observation~1. Conversely, the three judgment sets of \cref{cor:physics:eternalist} are
realizable but flagged: (i) $\{(\SUP(A),\emptyset,\bot)\}$ is realized by $M_@=\{N\}$ with
$N\models\neg A$; (ii) $\{(@,\emptyset,p_{\rm atm}=10^5),(c,\emptyset,p_{\rm atm}=0)\}$ and
(iii) $\{(c_1,\emptyset,Q=1),(c_2,\emptyset,Q=2)\}$ by one structure per context. Each
stripped set is unsatisfiable. The collapsed variant does not flag (i), but flags (ii) and
(iii).
\end{proof}

% -------------------------------------------------------------------
\subsection{Proofs for \S\ref{sec:physics:coherence}: blame localization and hygiene}

\begin{proof}[Proof of \cref{cor:physics:blame}]
(a) Let $w$ be a world at which $d$ is proper. The root is proper at every $w\in W_\rho$,
and $W_\rho\neq\emptyset$. If every Ax in the cone were true at $w$ and every Step, Imp and
Exp instance in it sound at $w$, \cref{thm:physics:soundness} (which uses only the items
of the derivation tree, i.e.\ the cone) would give $d\models_w\bot$, i.e.\
$\emptyset\in\fF_d(w)$. (b) A derivation of $d\cder\bot$ contains only the items in its tree,
and no rule of \cref{def:physics:calculus} takes $\ist$-claims, or a $\SUP$ context's
undischarged conclusions, as premises. (c) Let $\pi d$ be proper at $w$ with
$\pi d\models_w\omega_d$. Take a generator $\delta(F\times G)\cap\dom$ of $\fF_d(w)$. Since
$F\cap\tset{\omega_d}\in\fF_{\pi d}(w)$ and $\pi d$ is proper, there is
$\lambda\in F\cap\tset{\omega_d}$. Then $E_\lambda=\{\mu:\lambda[D{:=}\mu]\in\dom\}\in\fG$,
and since $\fG$ is proper, $G\cap E_\lambda\neq\emptyset$. For $\mu$ in it,
$\lambda[D{:=}\mu]\in\delta(F\times G)\cap\dom$. So every generator is nonempty, and
$d$ is proper by \cref{lem:physics:base}(a).
\end{proof}

\begin{proposition}[Propagation of $\bot$ along informative bridges]\src{T3 Cor 2.6, propagation clause}\label{prop:physics:propagation}
Let $c$ be a $\DEF$ context with an informative value-export bridge $\beta$ cited for $c$,
and suppose $c\cder\bot$, $\pi c\cder\sigma_\beta$ and $\pi c\cder\omega_c$ are derived.
Suppose $\hat R$ contains ex falso ($\bot\step\psi$) and the kernel of $\pi c$ is
refutation-complete for finite sets (it derives $\bot$ from any finite set of derived
sentences that has no model in $\fK$). Then $\pi c\cder\bot$ is derivable by adding finitely
many synthetic Exp instances. Iterating along an anchoring chain reaches the root or a
certified context. Moreover, fix a world $w$ at which $\pi c$ is proper:
\textup{(i)} if $c$ is proper at $w$, some item of $c$'s own derivation of $\bot$ is unsound
at $w$; \textup{(ii)} if $c$ is improper at $w$, all of $c$'s own items are sound at $w$, and
some item of the derivation of $\pi c\cder\omega_c$ is unsound at $w$; \textup{(iii)}
without the certificate premise in Exp, the unsound items in case (ii) are among the
synthetic Exp instances, i.e.\ the bridge applied to an ill-posed context.
\end{proposition}
\begin{proof}
Informativeness gives $t_1,\dots,t_k$ with $\{\sigma_\beta,\varepsilon_\beta(t_1),\dots,\varepsilon_\beta(t_k)\}$
unsatisfiable in $\fK$. Ex falso gives $c\cder Q=\bar t_i$ for each $i$, and $k$ Exp
instances give $\pi c\cder\varepsilon_\beta(t_i)$. Refutation completeness gives
$\pi c\cder\bot$. Iterate for $\pi c$ along the chain by which $c$ entered $\Anc$.

(i) is the contrapositive of \cref{thm:physics:soundness} in $c$. (ii) If $c$ is improper at
$w$, every set lies in $\fF_c(w)$, so every Step and Imp instance into or inside $c$ is
sound at $w$ and every Ax of $c$ is true there. By \cref{cor:physics:blame}(c),
$\pi c\not\models_w\omega_c$, so by \cref{thm:physics:soundness} the derivation of
$\pi c\cder\omega_c$ contains an item unsound at $w$. (iii) Without $\omega_c$, suppose the
derivation of $\sigma_\beta$ is sound at $w$, so $\pi c\models_w\sigma_\beta$. If every synthetic
instance were sound at $w$, then, as $c\models_w Q=\bar t_i$ trivially, $\pi c\models_w\varepsilon_\beta(t_i)$
for all $i$, and the finitely many truth sets would have empty intersection in the proper
filter $\fF_{\pi c}(w)$.
\end{proof}

\begin{proof}[Proof of \cref{thm:physics:hygiene}]
(a) Let $c'=\SUP(A)$ be the child of $c$ in which the reductio is carried out. By
\cref{lem:physics:base}(b), $\fF_{c'}(w)\supseteq\fF_c(w)$, and $\tset{A}\in\fF_{c'}(w)$. So the
premises $K_c$ (true in $c$) are true in $c'$, $A$ is true in $c'$, and every $c$-sound
instance has its truth set in $\fF_c(w)\subseteq\fF_{c'}(w)$, i.e.\ is $c'$-sound. By
induction along the derivation (filter closure, \cref{prop:physics:los}(a)) every line is
true in $c'$, so $c'\models_w\bot$, and \cref{prop:physics:sup} gives $c\models_w A\to\bot$,
i.e.\ $c\models_w\neg A$.

(b) A derivation of $\bot$ from $K_c$ is also a derivation of $\bot$ from $K_c\cup\{B\}$ for
every $B$, in particular for $B=A$ and $B=\neg A$. The set of refuted sentences is then all
sentences, whatever the world.

(c) From $A\to q$ and $\neg A\to q$, case analysis gives $q$, contradicting $\neg q$; so $K$
is unsatisfiable. The valuation $A=q=0$ satisfies $\{A\to q,\neg q\}$, and $A=1,q=0$
satisfies $\{\neg A\to q,\neg q\}$. $D_A$ derives $q$ from $A$ and $A\to q$ (modus ponens)
and then $\bot$ from $q,\neg q$; $D_{\neg A}$ is the same with $\neg A$. As a derivation
(used premises, steps, conclusion), $D_A$ is identical to the reductio of $A$ from the
consistent set $K'=\{A\to q,\neg q\}$, and that reductio is sound because $K'\cup\{A\}$ is
unsatisfiable. So a test that depends only on the derivation and accepts every sound
reductio from a consistent premise set accepts $D_A$; likewise $D_{\neg A}$ with
$K''=\{\neg A\to q,\neg q\}$. The accepted conclusions $\neg A$ and $\neg\neg A$ are jointly
unsatisfiable. (\texttt{hygiene.py} checks these facts by truth tables.)

(d) Suppose $N\models A$. Every leaf of the reductio is in $K_c\cup\{A\}$, hence true in $N$;
every step has true premises and an instance true in $N$, hence a true conclusion. So
$N\models\bot$, which is impossible; thus $N\models\neg A$. If both $A$ and $\neg A$ were
refuted by reductios whose instances $N$ satisfies, then $N\models\neg A$ and
$N\models\neg\neg A$. For the counterexample: if $c\models q$, the zero-premise instance
${}\step q$ is $c$-sound; with $K_c=\{\neg q\}$, every $\SUP(B)$ derives $q$ and $\neg q$,
hence $\bot$ by the sound instance $q,\neg q\step\bot$, so every $B$ is refuted; a model of
$\neg q$ exists but violates the instance ${}\step q$. For existence: if $c$ is proper at
$w$, $c\models_w K_c$ and the finitely many instances used are $c$-sound at $w$, then the
truth sets of the members of $K_c$ and of the instance conditionals
$\bigwedge\varphi_i\to\psi$ all lie in $\fF_c(w)$; their intersection lies in $\fF_c(w)$
and is nonempty, and any $\lambda$ in it gives $N=\fM_w(\lambda)$.

(e) By (b), each $A$ yields both labels; they are jointly unsatisfiable, a fixed function
of $K_c$, and the same in every world. If $K_c$ is consistent and $\hat R$ logically sound, a
model of $K_c$ satisfies every instance, so by (d) no $A$ and $\neg A$ are both refuted.
\end{proof}

% -------------------------------------------------------------------
\subsection{Proofs for \S\ref{sec:physics:export}}

\begin{proof}[Proof of \cref{thm:physics:nofree}]
In each case we exhibit $f$ meeting the hypotheses of \cref{lem:physics:invisible}.

(a) Let $f(\lambda)=e^{-1/\lambda^2}$ for $\lambda>0$ and $f(0)=0$. It is $C^\infty$ on
$[0,\Lambda]$ with all derivatives $0$ at $0$, and $f(\lambda^*)>0$. For $Q\in\mathcal C$ and
$A\in\R$, $Af\in C^\infty$, so $Q+Af\in\mathcal C$ by the closure hypothesis, and $Q+Af$ has
the same jet at $0$ as $Q$.

(b) Since $\lambda^*\notin\overline S$, some interval $(\lambda^*-r,\lambda^*+r)$ misses $S$.
Let $f$ be a $C^\infty$ bump supported in it with $f(\lambda^*)=1$. Then $Af\in C^\infty$,
$Q+Af\in\mathcal C$, and $(Q+Af)|_S=Q|_S$. The germ at $0$ is the special case
$S=[0,\lambda^*/2)$, since two functions that agree on $[0,\lambda^*/2)$ have the same germ.

(c) Let $f(\lambda)=\lambda^{N+1}$. It is entire, so $Q+Af$ is (the restriction of) an entire
function; its derivatives of order $\le N$ at $0$ vanish; and $f(\lambda^*)\neq0$.

(d) is (b) with the finite, hence closed, set $S=\{\lambda_1,\dots,\lambda_m\}$.

\emph{Weaker hypothesis.} If only $\mathcal C\supseteq C^\infty$, a rule sound on $\mathcal C$ is
sound on $C^\infty$, which satisfies $C^\infty+C^\infty\subseteq C^\infty$, so the rule
abstains on every $Q\in C^\infty$. Outside $C^\infty$ it need not: let
$\mathcal C=C^\infty\cup\{\sqrt\lambda\}$, $I$ the germ at $0$, and let the rule output
$[\sqrt{\lambda^*},\sqrt{\lambda^*}]$ when shown the germ of $\sqrt\lambda$ and abstain
otherwise. No smooth function has the germ of $\sqrt\lambda$ (it is not differentiable at
$0$ from the right), so the rule is sound on $\mathcal C$, and it does not abstain on
$\sqrt\lambda$.
\end{proof}

\begin{proof}[Proof of \cref{thm:physics:neighbourhood}]
(a) At every world compatible with the parent's derivation, $\lambda^*\in U$, so
$|Q(\fM(\lambda^*))-q|\le\epsilon$ by the frame theorem. (b) The negative half is
\cref{thm:physics:nofree}(b) (the germ is information on a set whose closure omits
$\lambda^*$); the positive half is (a) with $U$ and the stated rate. (c) The root filter is
the declared germ, so ``$Q\sim q$'' is a root claim; the decidability statement is
Gruntz's algorithm for limits of exp-log functions, modulo zero-equivalence of constants
\citep{gruntz1996computing}, which is undecidable for richer classes
\citep{richardson1968some}.
\end{proof}

\begin{proof}[Proof of \cref{thm:physics:gronwall}]
(a) Let $u(t)=|x(t)-y(t)|$. From
$x(t)-y(t)=x_0-y_0+\int_0^t\big(f(x)-f(y)-g(s,y)\big)ds$, the Lipschitz bound on $K$ and
$|g(s,y(s))|\le\eta$ give $u(t)\le u(0)+\int_0^t(Lu+\eta)\,ds=:v(t)$. Then $v$ is $C^1$ with
$v'=Lu+\eta\le Lv+\eta$, so $\frac{d}{dt}\big(e^{-Lt}(v+\eta/L)\big)=e^{-Lt}(v'-Lv-\eta)\le0$,
and $u\le v\le(u(0)+\eta/L)e^{Lt}-\eta/L=B(t)$.

(b) Let $t_1=\sup\{t\in[0,T]:|x(s)-y(s)|\le r\text{ for all }s\in[0,t]\}$. The set contains
$0$, because $|x_0-y_0|=B(0)\le B(T)<r$, and it is closed by continuity, so the supremum is
attained. On $[0,t_1]$ both $x(s)$ and $y(s)$ lie in $K_r$, and $|g(s,y(s))|\le\eta$ there;
so (a) applies on $[0,t_1]$ with $K=K_r$ and gives $|x-y|\le B(t)\le B(T)<r$. If $t_1<T$,
continuity gives $|x-y|<r$ on a neighbourhood of $t_1$, contradicting the definition of
$t_1$. Hence $t_1=T$.
\end{proof}

\begin{proof}[Proof of \cref{thm:physics:projectile}]
Write $v=(v_x,v_y)$, $v(0)=v_0(\cos\theta,\sin\theta)=(v_{0x},v_{0y})$, and
$\dot v=-g\hat y-\kappa v$; existence of the solution up to landing is the root model's
well-posedness. Let $T=\inf\{t>0:y(t)=0\}$; since $\dot y(0)=v_{0y}>0$, $y>0$ on $(0,T)$.

\emph{Speed bound.} On $[0,T]$,
$\frac{d}{dt}(\frac12|v|^2+gy)=v\cdot\dot v+gv_y=-\kappa|v|^2\le0$, and $y\ge0$; so
$|v|\le v_0$ and $|a_n|=\kappa|v|\le k|v|^2\le kv_0^2=xg$.

\emph{Flight time.} On $[0,T]$, $\ddot y=-g+a_{n,y}\in[-g(1+x),-g(1-x)]$, so
$v_{0y}t-\frac12g(1+x)t^2\le y(t)\le v_{0y}t-\frac12g(1-x)t^2$. The upper bound vanishes at
$T_{hi}=2v_{0y}/(g(1-x))$ and is negative after it; if $T>T_{hi}$ (or $T=\infty$), then
$y(T_{hi})\le0$ with $T_{hi}\in(0,T)$, a contradiction. So $T\le T_{hi}$. The lower bound is
positive on $(0,T_{lo})$ with $T_{lo}=2v_{0y}/(g(1+x))$; if $T<T_{lo}$ then $y(T)>0$, a
contradiction. So $T\ge T_{lo}$.

\emph{Horizontal motion.} $\dot v_x=-\kappa v_x$ gives $v_x=v_{0x}e^{-\int_0^t\kappa}>0$, so $X$
is increasing, and $\ddot X=-\kappa v_x\in[-\kappa|v|,0]\subseteq[-xg,0]$. Hence
$R=X(T)\le v_{0x}T\le v_{0x}T_{hi}=2v_{0x}v_{0y}/(g(1-x))=R_0/(1-x)$, using
$R_0=2v_{0x}v_{0y}/g$. And
$R\ge X(T_{lo})\ge v_{0x}T_{lo}-\frac12xgT_{lo}^2=\frac{R_0}{1+x}-\frac{2xv_{0y}^2}{g(1+x)^2}$;
since $R_0\tan\theta=2v_{0y}^2/g$, this is the stated lower bound.
\end{proof}

\begin{proof}[Proof of \cref{thm:physics:chains}]
(a) Triangle inequality. (b) The segment from $q$ to $Q$ lies in the convex box, on which
$F$ is differentiable; the mean value theorem gives $F(Q)-F(q)=\nabla F(\xi)\cdot(Q-q)$ for
some $\xi$ on the segment, and $|\nabla F(\xi)\cdot(Q-q)|\le\sum_jL_j\delta_j$. (c) If the
budget is at least $|F(q)|$, the certified interval
$[F(q)-\sum L_j\delta_j,F(q)+\sum L_j\delta_j]$ contains $0$. With $Q_1\in[99.8,100.2]$ and
$Q_2\in[99.7,100.1]$, $Q_1-Q_2\in[-0.3,0.5]$. For $F(q)=q^2$ on $[0.4,1.6]$, $L=3.2$ and the
budget is $1.92$, while $F([0.4,1.6])=[0.16,2.56]$.
\end{proof}

\begin{proof}[Proof of \cref{thm:physics:stipulation}]
(a) In a child with canonical stipulation $p_D=0$ (e.g.\ $k=0$), the smallness parameter
(e.g.\ $s=kv_0^2/g$) is $0$, so $c\cder s\le\eta$ is derivable soundly in the child, whatever
the world. An Exp rule that accepts it therefore exports $Q(\lambda^*)\in[Q(0)\pm\epsilon]$
from the value of $Q$ at the ideal point alone. This is an export rule whose information is a
single sample at $0\neq\lambda^*$, so by \cref{thm:physics:nofree}(d) it is unsound on any
class closed under adding smooth functions, for every $\epsilon$. The ping-pong numbers are
those of \cref{thm:physics:projectile} \status{computed}.

(b) With $k$ and $g$ both stipulated in the child, the child's range is
$R_0(g')=v_0^2\sin2\theta/g'$, which takes any positive value as $g'$ varies, while the
parent condition $kv_0^2/g\le\eta$ is evaluated at the actual $g$ and holds for the steel ball
($0.0276$). The export $R\in[R_0(g')\pm\epsilon]$ is false for $g'$ far from $g$.

(c) By \cref{thm:physics:soundness} and \cref{lem:physics:import}(a), it suffices that each
Exp instance is sound at every $w$. Let $p=\pi c$ and suppose $c\models_w Q=\bar v$,
$p\models_w\sigma_\beta$ and $p\models_w\omega_c$. By \cref{lem:physics:base}(a) there are
$F_1\in\fF_p(w)$ and $G\in\fG$ with $\delta(F_1\times G)\cap\dom\subseteq\tset{Q=\bar v}$. Let
$F=F_1\cap\tset{\sigma_\beta}\cap\tset{\omega_c}\in\fF_p(w)$ and fix $\lambda\in F$. Since
$\lambda\in\tset{\omega_c}$, $E_\lambda=\{\mu:\lambda[D{:=}\mu]\in\dom\}\in\fG$. For
$\mu\in G\cap E_\lambda$, $\lambda[D{:=}\mu]\in\delta(F_1\times G)\cap\dom$, so
$Q(\fM_w(\lambda[D{:=}\mu]))=v$. Hence
$\{\mu:\lambda[D{:=}\mu]\in\dom,\ Q(\fM_w(\lambda[D{:=}\mu]))=v\}\supseteq G\cap E_\lambda\in\fG$,
and the schema's antecedent holds at $\lambda$; as $\fM_w(\lambda)\models\sigma_\beta$, the
schema gives $\fM_w(\lambda)\models\varepsilon_\beta(v)$. So
$F\subseteq\tset{\varepsilon_\beta(v)}$, i.e.\ $p\models_w\varepsilon_\beta(v)$.

\emph{The rope.} $\dom=\{m_r>0\}$, $D=\{m_r\}$, limit semantics. The child's generator
$\{\lambda^*[m_r{:=}0]\}\cap\dom$ is empty, so the child is improper and proves $a=0$ by
(vacuously) sound steps. The schema's antecedent requires $\lambda[m_r{:=}0]\in\dom$, which
never holds, so the schema is a theorem of the frame. The parent proves $m_r\le\eta$ when
the actual $m_r\le\eta$, and Exp without $\omega_c$ would export $|a|\le\epsilon$, whereas
$|a|=|\Phi|/m_r\ge|\Phi|/\eta>\epsilon$ for $\eta<|\Phi|/\epsilon$. With the certificate
requirement, $\tset{\omega_c}$ must be empty (no $\lambda$ has $\lambda[m_r{:=}0]\in\dom$),
so the proper root cannot soundly derive $\omega_c$.
\end{proof}

\begin{lemma}[Asymptotic schemas below a principal parent]\src{T3 Remark after Thm 3.9}\label{lem:physics:asymptotic}
Read the child judgment ``$Q\simeq v$'' as: $c\models_w|Q-\bar v|<\bar r$ for every rational
$r>0$. Call $\beta$ an \emph{asymptotic schema} for the declaration $(D,\lambda^\circ,\fG)$ if,
for all $w$, all $\lambda\in\dom\fM_w$ with $\fM_w(\lambda)\models\sigma_\beta$ and
$E_\lambda=\{\mu:\lambda[D{:=}\mu]\in\dom\fM_w\}\in\fG$, and all $v$: if
$Q(\fM_w(\lambda[D{:=}\mu]))\to v$ along the trace of $\fG$ on $E_\lambda$, then
$\fM_w(\lambda)\models\varepsilon_\beta(v)$. If $\pi c$ is principal at $w$ (the root, or a
limit-semantics context below principal ancestors), the Exp instance from $c\cder Q\simeq v$,
$\pi c\cder\sigma_\beta$ and $\pi c\cder\omega_c$ to $\pi c\cder\varepsilon_\beta(v)$ is sound
at $w$.
\end{lemma}
\begin{proof}
Let $\pi c$ be principal at $\lambda$ (if it is improper, soundness is vacuous). Then
$\fM_w(\lambda)\models\sigma_\beta$ and $\lambda\in\tset{\omega_c}$, so $E_\lambda\in\fG$. By
\cref{lem:physics:base}(c), $c\models_w|Q-\bar v|<\bar r$ iff there is $G\in\fG$ with
$|Q(\fM_w(\lambda[D{:=}\mu]))-v|<r$ for all $\mu\in G\cap E_\lambda$. Holding for every
rational $r>0$, this is convergence to $v$ along the trace of $\fG$ on $E_\lambda$. The
schema gives $\fM_w(\lambda)\models\varepsilon_\beta(v)$, i.e.\ $\pi c\models_w\varepsilon_\beta(v)$.
Without $\omega_c$ the argument fails at $E_\lambda\in\fG$: an improper child proves $Q\simeq v$
for every $v$. Under a non-principal parent the convergence would have to be uniform on a
parent filter set, which is an additional side condition.
\end{proof}

% -------------------------------------------------------------------
\subsection{Proofs for \S\ref{sec:physics:learning}}

\begin{proof}[Proof of \cref{thm:physics:conformal}]
(a) If $k>n$ then $\hat\epsilon=\infty$ and there is nothing to prove; note $k\ge1$ for
$\alpha<1$. Call index $i\in\{1,\dots,n+1\}$ \emph{high} if $e_i$ is strictly greater than at
least $k$ of the other $n$ values. If $e_{n+1}>\hat\epsilon$, the $k$-th smallest of
$e_1,\dots,e_n$, then $e_{n+1}$ exceeds at least $k$ of them, so $n+1$ is high. In any
realization at most $n+1-k$ indices are high (in sorted order, only positions $k+1,\dots,n+1$
can have $k$ strictly smaller values). By exchangeability every index is high with the same
probability $p$, so $(n+1)p=\E[\#\text{high}]\le n+1-k$, and
$\Prob(e_{n+1}>\hat\epsilon)\le p\le(n+1-k)/(n+1)\le\alpha$ because $k\ge(n+1)(1-\alpha)$.
Ties are handled by the strict inequality. (b) If the chosen instance has $e>\hat\epsilon$,
it is not covered, and the solver can choose such an instance whenever one exists.
\end{proof}

\begin{proof}[Proof of \cref{thm:physics:noreg}]
We take an exact oracle (or any oracle whose answer is a fixed function of the value at the
queried point); this only strengthens the certifier.

\emph{Deterministic.} Fix $e\in\mathcal C$, let $\pi_1,\dots,\pi_q$ be the queries made on $e$
and $V$ the output. Suppose $\pi\in V$ is not queried. Let $f\ge0$ be a continuous bump with
$f(\pi)=1$ whose support avoids the finite set $\{\pi_1,\dots,\pi_q\}$, and let $A>0$. On
$e'=e+Af\in\mathcal C$ the first query is $\pi_1$, with the same answer since
$f(\pi_1)=0$; by induction the whole adaptive query sequence and the output coincide, so
$\pi\in V(e')$. Soundness on $e'$ requires $e(\pi)+A\le\tau$ for all $A>0$, which is
impossible.

\emph{Randomized.} Suppose that on some $e\in\mathcal C$ the event that $V$ contains an
unqueried point has positive probability (we assume the events below are measurable). For
each outcome in this event there is an unqueried $\pi\in V$ and a ball $B$ from a fixed
countable base (rational centres and radii) with $\pi\in B$ and $\overline B$ disjoint from the
finitely many queried points. Some fixed $B$ occurs with positive probability. Let
$f(\pi')=\mathrm{dist}(\pi',\Pi\setminus B)$, a continuous bump, positive exactly on $B$. For
outcomes in the event for $B$, the transcript on $e+Af$ equals that on $e$, so $V$ still
contains a point $\pi\in B$ with $f(\pi)>0$, and $e(\pi)+Af(\pi)>\tau$ once $A$ is large
enough for that outcome. By monotone convergence, some fixed $A$ makes the certifier unsound on
$e+Af\in\mathcal C$ with positive probability.
\end{proof}

\begin{proof}[Proof of \cref{thm:physics:lipschitz}]
(a) If $\pi\in V$, then $\|\pi-\pi_j\|_\infty\le(\tau-e_{hi}(\pi_j))/L$ for some centre
$\pi_j$, so $e(\pi)\le e(\pi_j)+L\|\pi-\pi_j\|_\infty\le e_{hi}(\pi_j)+(\tau-e_{hi}(\pi_j))=\tau$.

(b) Let $e(\pi)\le\tau-\gamma$, and let $\pi_j$ be the centre of a cell containing $\pi$, so
$\|\pi-\pi_j\|_\infty\le h/2$. Then $e_{hi}(\pi_j)\le e(\pi_j)+\eta\le e(\pi)+Lh/2+\eta\le
\tau-\gamma+Lh/2+\eta$, so the radius at $\pi_j$ is at least
$(\gamma-\eta)/L-h/2\ge h-h/2=h/2\ge\|\pi-\pi_j\|_\infty$, and $\pi\in V$. With
$h=(\gamma-\eta)/L$ there are $\lceil1/h\rceil^d=\lceil L/(\gamma-\eta)\rceil^d$ centres.

(c) Let $e_0\equiv\tau-\gamma$; completeness forces $V(e_0)=\Pi$. Put $\rho=\gamma'/L$ and pack
$M=\lfloor1/(2\rho)\rfloor^d$ disjoint open $\|\cdot\|_\infty$-balls of radius $\rho$ in $\Pi$.
With fewer than $M$ queries on $e_0$, some ball $B(\pi_0,\rho)$ contains no query. Define
$e_1=e_0+\max(0,\gamma'-L\|\pi-\pi_0\|_\infty)$. It is $L$-Lipschitz, equals $e_0$ outside the
open ball, and $e_1(\pi_0)=\tau-\gamma+\gamma'>\tau$. All answers on $e_1$ agree with those
on $e_0$, so the deterministic certifier outputs $V(e_1)=\Pi\ni\pi_0$, which is unsound.
\end{proof}

\begin{proof}[Proof of \cref{thm:physics:monotone}]
(a) $lo$ changes only to accepted points, so if $V=[0,lo]$ then $e_{hi}(lo)\le\tau$ was
observed, and for $\pi\le lo$, $e(\pi)\le e(lo)\le e_{hi}(lo)\le\tau$.

(b) Each query halves $hi-lo$, so $hi-lo=2^{-k}$ at the end. With an exact oracle a rejection
at $hi$ means $e(hi)>\tau$, so by monotonicity $\{e\le\tau\}\subseteq[0,hi)$; if nothing was
rejected, $hi=1$ and $\{e\le\tau\}\subseteq[0,1]$. If $V\neq\emptyset$, the missed part lies
in $(lo,hi]$, of length $2^{-k}$. If $V=\emptyset$, every query was rejected, so $hi=2^{-k}$
and the missed part lies in $[0,hi)$. The endpoint variant spends two queries before
bisecting, leaving $k-2$ halvings.

(c) Let $\tau>0$ and $e_t=2\tau\mathbf 1[\pi\ge t]$, non-decreasing with $\{e_t\le\tau\}=[0,t)$.
Exact answers on $e_t$ take only the values $0$ and $2\tau$, so a deterministic certifier with
$q$ queries has at most $2^q$ transcripts and hence at most $2^q$ outputs. Choose
$M=\lceil1/h\rceil$ thresholds in $(0,1]$ pairwise more than $h$ apart (possible since
$(M-1)h<1$). If $2^q<M$, two thresholds $t<t'$ share a transcript and hence an output $V$.
Soundness on $e_t$ gives $V\subseteq[0,t)$, so on $e_{t'}$ the missed part contains $[t,t')$, of
length more than $h$. Hence $2^q\ge M\ge1/h$, i.e.\ $q\ge\lceil\log_2(1/h)\rceil$.

(d) \emph{Sketch.} On a grid of side $h$, the cells meeting a hyperplane
$\pi_1+\dots+\pi_d=\mathrm{const}$ form an antichain of size $\Theta(h^{1-d})$ for the
coordinatewise order; down-sets can include or exclude each of them independently, so each
must be probed. A staircase walk along the boundary of the down-set achieves the bound.
\end{proof}

% -------------------------------------------------------------------
\subsection{Proofs for \S\ref{sec:physics:checker}}

\begin{proof}[Proof of \cref{thm:physics:relative}]
Fix $w\in W_\rho$. We check that every Ax in the accepted derivation is true at $w$ and every
Step, Imp and Exp instance is sound at $w$; then \cref{thm:physics:soundness} gives
$@\models_w Q\in[q-\delta,q+\delta]$, and since $\fF_@(w)$ is principal at $\lambda^*_w$,
$Q^{\fM_w(\lambda^*_w)}\in[q-\delta,q+\delta]$.
\begin{itemize}
\item \emph{Step.} By (TB1) the instance conditional $\bigwedge\varphi_i\to\psi$ (and every
  cited law) is valid in every $\fM_w(\lambda)$, $\lambda\in\dom$; its truth set is $\dom$,
  which lies in every filter on $\dom$. So the instance is $c$-sound for every context $c$.
\item \emph{Root axioms} lie in $\Ax_\rho$, which is true at $\fM_w(\lambda^*_w)$ by the
  definition of a reading, hence in the principal root filter.
\item \emph{$\SUP$ axioms:} $\tset{A}\in\fF_{\SUP(A)}(w)$ by definition.
\item \emph{$\DEF$ axioms} are canonical (Ax legality). Under limit semantics
  $\{\lambda^\circ\}\in\fG$, and the generator $\delta(F\times\{\lambda^\circ\})\cap\dom$ lies
  in $\tset{p_D=\lambda^\circ}$. Under germ semantics, for definable $U\in\fG$ the generator
  $\delta(F\times U)\cap\dom$ lies in $\tset{p_D\in U}$. This check is what blocks
  \cref{thm:physics:stipulation}(b).
\item \emph{Imports into $\DEF$ contexts} are catalogue laws (valid on $\dom$ by (TB1)) or
  undeformed data, hence $D$-stable at $w$, hence sound by \cref{lem:physics:import}(a).
  Imports into $\SUP$ children are sound because $\fF_{\SUP(A)}(w)\supseteq\fF_p(w)$.
  $\AUX$ bridges are kernel-checked theorems, sound by (TB1).
\item \emph{Discharges} are sound by \cref{prop:physics:sup}.
\item \emph{Exports.} Schema indexing ensures that the cited $\beta$ is, by (TB2), a theorem
  for the context's actual declaration. In the pointwise form, the core of the proof of
  \cref{thm:physics:stipulation}(c) shows that the instance is sound at $w$; in the
  asymptotic form, cited only below a principal parent, \cref{lem:physics:asymptotic} does.
  Validated enclosures used in side conditions are correct by (TB3). The premises
  $\sigma_\beta$ and $\omega_c$ are derived in the parent, so in the induction of
  \cref{thm:physics:soundness} they are true there at $w$.
\end{itemize}
Nothing in this argument depends on how the derivation was produced, so it holds for every
solver. For a germ root $d$ with trusted declaration, the same argument yields
$d\models_w\chi$. If $W_{\rho^*}\subseteq W_\rho$, the conclusion holds at every
$w\in W_{\rho^*}$, i.e.\ it is supertrue for the intended problem; if the actual system is a
world in $W_{\rho^*}$, the conclusion is true of it.
\end{proof}

\begin{proof}[Proof of \cref{thm:physics:superval}]
(a) Supertruth of $Q\in[q\pm\delta]$ means $|Q^w-q|\le\delta$ for every $w\in W_\rho$, i.e.\
$\sup_w|Q^w-q|\le\delta$. If some $q$ satisfies this, all $Q^w$ lie in an interval of length
$2\delta$, so $\osc Q\le2\delta$. Conversely, if $\osc Q\le2\delta$ (and $W_\rho\neq\emptyset$),
the midrange $q=\frac12(\sup Q^w+\inf Q^w)$ has $|Q^w-q|\le\frac12\osc Q\le\delta$. (b) is
\cref{thm:physics:relative} applied to root derivations. (c) By
\cref{thm:physics:judgment}(a), a checker sound for every $\rho^*\in\mathcal R$ accepts only
claims holding on $\bigcup_{\rho\in\mathcal R}W_\rho$; and accepting exactly the claims
supertrue on that union is sound for every $\rho^*\in\mathcal R$, since
$W_{\rho^*}\subseteq\bigcup_{\rho\in\mathcal R}W_\rho$.
\end{proof}

% -------------------------------------------------------------------
\subsection{Proofs for \S\ref{sec:physics:example}}

\begin{proof}[Proof of \cref{prop:physics:legexact}]
(a) The incoming direction is $u=(\sin\alpha\,e(0),-\cos\alpha)$. Reflection in the vertical
surface with outward normal $n=(e(\vartheta),0)$ keeps the vertical component and maps the
horizontal one to $\sin\alpha\,(e(0)-2\cos\vartheta\,e(\vartheta))=\sin\alpha\,e(2\vartheta-\pi)$;
the surface is lit where $u\cdot n<0$, i.e.\ $\cos\vartheta<0$. Descending a height $h$ takes
horizontal travel $h\tan\alpha=s$, so $P=a\,e(\vartheta)+s\,e(\psi)$. For the second form,
$e(\vartheta)\cdot e(\psi)=\cos(\pi-\vartheta)=-\cos\vartheta=\sin\frac\psi2$ and
$e(\vartheta)\cdot e(\psi+\frac\pi2)=\sin\vartheta=\cos\frac\psi2$.

(b) A direct computation (checked in SymPy) gives
$\det\partial P/\partial(\vartheta,s)=a\cos\vartheta-2s$, which is negative on the lit side, so
$|\det|=2s-a\cos\vartheta=2s+ac$. For injectivity, note first that along one reflected ray
the coefficient $s+ac$ of $e(\psi)$ is strictly increasing in $s$, and distinct $\vartheta$
give distinct $\psi$. So it suffices that two distinct reflected rays never meet. By (a) the
ray with parameter $\psi$ is $\{t\,e(\psi)+a\cos\frac\psi2\,e(\psi+\frac\pi2):t\ge a\sin\frac\psi2\}$.
Take $\psi_1=2x<\psi_2=2y$ with $x,y\in(0,\pi)$ and put $d=y-x\in(0,\pi)$. If $\sin2d\neq0$,
solving the $2\times2$ linear system for the intersection of the two lines gives
\[
t_1=\frac{a(\cos x\cos2d-\cos y)}{\sin2d},\qquad t_2=\frac{a(\cos x-\cos y\cos2d)}{\sin2d}
\]
(checked in SymPy), and the rays meet only if $t_1\ge a\sin x$ and $t_2\ge a\sin y$.
\emph{Case $\sin2d>0$, i.e.\ $d<\pi/2$.} $t_2\ge a\sin y$ is $\cos x\ge\cos(y-2d)=\cos(d-x)$;
as $x,|d-x|\in[0,\pi)$ and $\cos$ decreases there, $x\le|d-x|$, which forces $d\ge2x$.
$t_1\ge a\sin x$ is $\cos(x+2d)\ge\cos(x+d)$; with $u=x+d\in(0,\pi)$ and $u<x+2d<2\pi$, this
forces $x+2d>\pi$ and $2\pi-(x+2d)\le u$, i.e.\ $2x+3d\ge2\pi$. Then $4d\ge2x+3d\ge2\pi$,
contradicting $d<\pi/2$. \emph{Case $\sin2d<0$, i.e.\ $d>\pi/2$.} The inequalities reverse:
$\cos x\le\cos(d-x)$ gives $x\ge|d-x|$, hence $d\le2x$; and $\cos(x+2d)\le\cos(x+d)$, where
$u=x+d\in(\pi/2,\pi)$ and $x+2d\in(\pi,2\pi)$, gives $2\pi-(x+2d)\ge u$, i.e.\ $2x+3d\le2\pi$.
Then $4d\le2x+3d\le2\pi$, contradicting $d>\pi/2$. \emph{Case $d=\pi/2$.} Then
$e(2y)=-e(2x)$, the lines are parallel, and they coincide only if $a\cos x=-a\cos y=a\sin x$,
i.e.\ $x=\pi/4$; then $t_1e(2x)$ and $-t_2e(2x)$ with $t_1,t_2>0$ point apart, so the rays
do not meet. (A numerical search over $2\times10^6$ random pairs also found no forward
intersection.)

(c) The irradiance on a surface element is $I\sin\alpha\,|\cos\vartheta|$ and the element has
area $a\,d\vartheta\,dh$ with $dh=ds/\tan\alpha$. So the reflected flux through
$d\vartheta\,ds$ is $I\cos\alpha\,a|\cos\vartheta|\,d\vartheta\,ds$; dividing by the floor area
$(2s+ac)\,d\vartheta\,ds$ and using $|\cos\vartheta|=c$, $I\cos\alpha=I_0$ gives
$E=I_0\,ac/(2s+ac)$. Floor points outside the image receive no reflected light, and points
whose preimage has $s>L\tan\alpha$ would need reflection above the lit height. The total flux
is $\int_0^L\!\!\int_{\pi/2}^{3\pi/2}I\sin\alpha\,|\cos\vartheta|\,a\,d\vartheta\,dh=2aLI\sin\alpha$,
the flux intercepted by the lit side.

(d) Neither $P(\vartheta,s)$ nor $E$ as a function of $(\vartheta,s,I_0)$ contains $\alpha$; only
the constraint $s\le L\tan\alpha$ does.
\end{proof}

\begin{proof}[Proof of \cref{prop:physics:thinleg}]
Since $\epsilon<2/\sqrt5<1$, $r>a$.

\emph{Step 0: the point is in the image and lit.} On the ray with parameter $\psi$, the point at
distance $r$ from the axis is $u\,e(\psi)+a\cos\frac\psi2\,e(\psi+\frac\pi2)$ with
$u=\sqrt{r^2-a^2\cos^2\frac\psi2}>0$, and it lies on the forward ray ($s=u-ac\ge0$) because
$u^2-a^2c^2=r^2-a^2\ge0$. Its polar angle is $\varphi(\psi)=\psi+\delta(\psi)$ with
$\tan\delta=a\cos\frac\psi2/u$, continuous in $\psi$, tending to $\arcsin\epsilon$ as
$\psi\to0^+$ and to $2\pi-\arcsin\epsilon$ as $\psi\to2\pi^-$. By the intermediate value
theorem every $\varphi\in(\arcsin\epsilon,2\pi-\arcsin\epsilon)$ is attained. And $\sigma>m$
forces $\varphi$ into this interval: for $\varphi\in[0,\arcsin\epsilon]$,
$\sigma=\sin\frac\varphi2\le\frac\varphi2\le m$, and symmetrically near $2\pi$. So the point
has a preimage, unique by \cref{prop:physics:legexact}(b), with $s\le u\le r\le L\tan\alpha$,
so it is lit.

\emph{Step 1.} With $u=s+ac\ge0$, \cref{prop:physics:legexact}(a) gives
$r^2=u^2+a^2\cos^2\frac\psi2$, so $u\in[r\sqrt{1-\epsilon^2},r]$.

\emph{Step 2.} $\sin\delta=a\cos\frac\psi2/r$, so $|\sin\delta|\le\epsilon$ and, as
$\delta\in(-\pi/2,\pi/2)$, $|\delta|\le\arcsin\epsilon$. Since $c=\sin\frac\psi2\ge0$ and
$\sin$ is $1$-Lipschitz, $|\sigma-c|\le|\sin\frac\varphi2-\sin\frac\psi2|\le|\delta|/2\le m$.

\emph{Step 3.} $E/E_0=\frac{ac}{2s+ac}\cdot\frac{2r}{a\sigma}=\frac c\sigma\cdot\frac{2r}{2u-ac}$.
The second factor is at least $1$ because $2u-ac\le2r$, and at most
$1/(\sqrt{1-\epsilon^2}-\epsilon/2)$ because $2u-ac\ge2r\sqrt{1-\epsilon^2}-a$; the denominator
is positive exactly when $\epsilon<2/\sqrt5$.

\emph{Step 4.} By Step 2, $c/\sigma\in[1-m/\sigma,1+m/\sigma]$, and $1-m/\sigma>0$. Multiplying the
two ranges of nonnegative factors gives the claim.
\end{proof}

\begin{proof}[Proof of \cref{prop:physics:finitesun}]
\emph{Domain.} Each direction $d$ in the solar disc has zenith angle
$\alpha_d\ge\alpha-\delta_s\ge\alpha_{\min}-\delta_s$. The preimage of a floor point at distance
$r$ has $s\le u\le r$ (Step 0 above), and $r\le L\tan(\alpha_{\min}-\delta_s)\le L\tan\alpha_d$,
so the point is lit by every direction, for every admissible $\alpha$.

\emph{Convex combination.} By \cref{prop:physics:legexact}(d) and rotational symmetry, direction
$d$ with azimuth offset $\beta_d$ contributes $G(r,\varphi-\beta_d)\cos\alpha_d\,dI(d)$, where
$G=E/I_0$ is the $\alpha$-free geometric factor. The direct illuminance is
$I_0=\int\cos\alpha_d\,dI(d)$, so $E_\odot/I_0$ is a convex combination of the values
$G(r,\varphi-\beta)$ over the azimuth spread, for any radiance profile over the disc.

\emph{Azimuth spread.} In the right spherical triangle formed by the zenith, the sun's centre
and a tangent point of the disc, Napier's rule gives $\sin\beta_{\max}=\sin\delta_s/\sin\alpha$.
Since $\alpha\in[\alpha_{\min},\pi/2]$ and $\delta_s<\alpha_{\min}$, this is at most
$\sin\delta_s/\sin\alpha_{\min}$, so $|\beta|\le\Delta$.

\emph{Combination.} Let $\sigma'=|\sin\frac{\varphi-\beta}2|$. As $\sin$ is $1$-Lipschitz,
$|\sigma'-\sigma|\le|\beta|/2\le\Delta/2$, so $\sigma'\ge\sigma-\Delta/2>m$ and
\cref{prop:physics:thinleg} applies at $\varphi-\beta$:
$G(r,\varphi-\beta)\in[\mathrm{lo}(\epsilon,\sigma'),\mathrm{hi}(\epsilon,\sigma')]\cdot E_0(\varphi-\beta)/I_0$.
Since $\mathrm{lo}$ increases and $\mathrm{hi}$ decreases in $\sigma$, these factors lie in
$[\mathrm{lo}(\epsilon,\sigma-\frac\Delta2),\mathrm{hi}(\epsilon,\sigma-\frac\Delta2)]$, and
$E_0(\varphi-\beta)/E_0(\varphi)=\sigma'/\sigma\in[1-\frac\Delta{2\sigma},1+\frac\Delta{2\sigma}]$.
All factors are nonnegative ($\sigma>\Delta/2+m$), so multiplying and averaging over the convex
combination gives the claim.
\end{proof}
